% !PRINTED 49
\weberpage{49}

If conversely one of the characters $\chi = \xi$ of the group $S$
has the same value for all elements of the group $T$, then this
can only be the value 1, since the unit element of $S$ certainly
occurs in $T$ (\S.\,11, 2.). If then $A_i$ is any element
of $T_i$, and $A$ runs through the whole group $T$, then
$AA_i$ runs through the subsidiary group $T_i$. But then
$\xi(A) = 1$ and consequently
\weq{5}{\xi(AA_i) = \xi(A_i),}
i.\,e. $\xi(A)$ has one and the same value for all elements of a
subsidiary group $T_i$, and $\xi(A)$ can therefore be regarded as a
function of $T_i$ and denoted by $\xi(T_i)$.

Since now, if $A_i$ occurs in $T_i$ and $A_k$ in $T_k$, $A_iA_k$
is contained in $T_iT_k$, it follows that
\weq{6}{\xi(T_i)\xi(T_k) = \xi(T_iT_k),}
i.\,e. $\xi(T_i)$ is contained among the characters of the group
$S/T$. There are therefore exactly $j$ such characters $\xi(A)$,
which are defined by the property that they have the value 1 for
every element in $T$. These $j$ characters $\xi$ form a group under
the composition of the characters; for if $\xi_1(A) = 1$,
$\xi_2(A) = 1$, then also $\xi_1\xi_2(A) = 1$ [\S.\,11, (18)].
Since by (4) the functions $\xi$ can also be regarded as the
characters of the group $S/T$, by \S.\,11, 5. the group of the
$\xi$ is isomorphic with the group $S/T$.

With this, then, the following theorem is proved:

7.\ \gesperrt{If an Abelian group $S$ has a divisor $T$ of
degree $t$ and of index $j$, then among the characters of
$S$ there are exactly $j$ and no more that have the value 1
for all elements of $T$, while for every element of $S$ not
contained in $T$ at least one of them is different from 1. These
characters form a group isomorphic with $S/T$.}

We shall call these characters \gesperrt{belonging to the group $T$}.

Since every character $\chi_B$ corresponds to a definite element $B$ of $S$,
it follows also that, when $\chi_B$ runs through the group of the
characters belonging to $T$, $B$, by the formula \S.\,11, (18),
will run through a group which is of degree $j$ and isomorphic with the
group of the $\chi_B$ and hence also with the group $S/T$. This
%CONT>>
